\section{Conclusion and Future Work}\label{sec:conclusion}

Across all layers of GPT-2 Small, three SAE architectures reconstruct activations well by conventional measures, but
reconstruction quality predicts behavioural damage rather than information preservation, frozen-probe ``truth loss''
is largely perturbation sensitivity, PCA at equal reconstruction error preserves at least as much task information
as the SAEs (while using more dimensions than the SAEs activate per token), and at mid and late
layers the SAE reconstruction carries only about two-thirds of a causal effect located in the original model.
Short fine-tuning improves the faithfulness level it optimises (task reconstruction, or output agreement with an
output-KL term) but not truth information as seen by a re-trained reader, and the behavioural gains of ReLU and
JumpReLU at depth come with a sparsity cost. On the persona task the SAEs lose information that re-trained readers
cannot recover and distort the honest-minus-dishonest difference, and neither fine-tuning method creditably improves
persona representation or the causal effect carried at layer 8.

\paragraph{Improving the SAEs.} Training our own SAEs on 1000$\times$ more data at 4$\times$ the width (R2,
Section~\ref{sec:improve}) changes several of these conclusions and leaves others intact.
\begin{itemize}
\item \textbf{Confirmed on held-out data:} reconstruction error falls by 31--74\,\% and splice damage by
  35--79\,\%.
\item \textbf{Confirmed on held-out RepE-LM data (H1-TEST):} persona pair-direction fidelity improves in all 9
  layer--comparison cells (+0.035 to +0.141 cosine), and for TopK the persona effect carried at layer 8 rises from
  0.602 to 0.729 (Holm-adjusted $p=0.021$).
\item \textbf{Not confirmed on held-out data:} the layer-8 causal gain for BatchTopK against JumpReLU (+0.072, all 9
  seed pairs better, Holm-adjusted $p=0.058$), and the persona re-fit-gap reduction as a hypothesis (TopK passes,
  $-0.048$, Holm-adjusted $p=0.046$; BatchTopK against JumpReLU is inconclusive). A test-only neutral control
  exceeded its floor at layers 8 and 11 (0.561, 0.551 against 0.55).
\item \textbf{Unchanged:}
  \begin{itemize}
  \item SAEs still do not beat PCA at equal error on task information;
  \item truth information shows no confirmed gain, within a design that cannot detect changes smaller than about
    0.05--0.09 AUROC;
  \item a supervised rank-4 subspace carries more of the persona effect at layer 8 than the full SAE reconstruction.
  \end{itemize}
\item \textbf{KL fine-tune at fixed L0:} it adds a confirmed 16--17\,\% behavioural gain, but reduces causal
  preservation at layer 11. Faithfulness levels must therefore be optimised and reported separately.
\item \textbf{External SAE:} the same conclusions hold for a public SAE trained on 300M tokens.
\item \textbf{The PCA comparison depends on the constraint} (\evlabel{validation}). At an equal 128-dimensional
  budget, trained SAEs are better than PCA under the study's comparison criterion in 3 of 20 GoT comparisons and not
  classified worse in the remaining 17; with PCA matched to task-specific SAE activity, better in 8 of 20 and not
  classified worse in 12 (no equivalence or non-inferiority is claimed).
\end{itemize}
Except where marked \evlabel{test}, results are \evlabel{validation} findings.

\paragraph{Future work.}
\begin{itemize}
\item Gemma-2-2B, using Gemma Scope SAEs as the external anchor and, where feasible, our own data-scaled TopK SAEs at
  one middle layer. The ladder transfers directly, with the KL fine-tune only as a behaviour-only add-on.
\item Calibrating neutral frames on training data before any future test-only template set, and more seeds for the
  BatchTopK-against-JumpReLU causal comparison (in a new, separately pre-registered test set).
\item Capacity-matched PCA for the persona task, and a pre-registered test-split confirmation of PCA-CM.
\item Task-level sparsity matching.
\item Supervised-subspace (DAS) skylines as a standard reference for causal SAE evaluations.
\end{itemize}
